%!TEX root = ../main/main.tex


\item \text{[2 puntos]} En esta pregunta usted deberá programar una clase que represente a un participante de un conjunto de personas que generan firmas de anillos trazables. También una clase que revise dichas firmas y que verifique si el participante que emitió la firma ya había emitido otra firma anteriormente. Utilizaremos siempre SHA256 y grupos multiplicativos de enteros módulo primos ($\mathbb{Z}_p^*$). En particular:

\begin{itemize}
        \item Para transformar un número \texttt{n} a un string, simplemente utilizaremos \texttt{str(n)}.
        \item Para transformar el hash de un string \texttt{s} a un elemento del grupo $\mathbb{Z}_p^*$ utilizaremos \texttt{int.from\_bytes(sha256(s.encode())) \% p}.
\end{itemize}

La clase que representa a un participante debe tener la siguiente interfaz:

\begin{python}
class TraceableRingSignatureParticipant:
        def __init__(self, g: int, q: int, p: int) -> None:
                # Verify that p is prime, raise an AssertionError otherwise.
                # Verify that q is prime in Zp*, raise an AssertionError otherwise.
                # Verify that the order of the subgroup generated by g is q, raise an AssertionError otherwise.
                # Randomly generate the participant's secret key x in {1..q}

        def get_public_key(self) -> int:
                # Returns the participant's public key (g^x mod p)

        def generate_traceable_ring_signature(self, public_keys: list[int], tag: str, message: str) -> Tuple[list[int], int, int]:
                # Compute T := (SHA256(p_1 || ... || p_n || tag)^r mod p), where [p_1, ..., p_n] is the list public_keys of public keys.
                # Verify that the order of the subgroup generated by T is q, raise an AssertionError otherwise.
                # Generate a traceable ring signature for the given message, tag and list of public keys.
                # The output looks like ([s1, ..., sn], c1, f).

\end{python}

La clase que representa a un verificador debe tener la siguiente interfaz:

\begin{python}
class Verifier:
        def __init__(self, g: int, q: int, p: int) -> None:
                # Verify that p is prime, raise an AssertionError otherwise.
                # Verify that q is prime in Zp*, raise an AssertionError otherwise.
                # Verify that the order of the subgroup generated by g is q, raise an AssertionError otherwise.

        def verify_signature(self, public_keys: list[int], message: str, tag: str, signature: Tuple[list[int], int, int]) -> bool, bool:
                # Verify if the signature is a valid traceable ring signature for the given public keys, message and tag.
                # The first component of the output is True if the signature is valid.
                # The second component is True if this is the first signature of this participant seen by the verifier.
                # Note the verifier needs to store some state to compute the second component.
\end{python}

\medskip

\paragraph{Corrección.}
La corrección de esta pregunta va a estar basada en el siguiente conjunto de tests.

\begin{itemize}
\item \text{[0.25 puntos]} Al inicializar un \texttt{TraceableRingSignatureParticipant} con un valor de $p$ que no es primo, se debe obtener una excepción.

\item \text{[0.25 puntos]} Al inicializar un \texttt{TraceableRingSignatureParticipant} con un valor de $p$ primo y un valor de $q$ que no es primo, se debe obtener una excepción.

\item \text{[0.25 puntos]} Al inicializar un \texttt{TraceableRingSignatureParticipant} con valores primos de $p$ y $q$, y con un elemento $g \in \mathbb{Z}_p^*$ tal que $|\langle g \rangle| \neq q$, se debe obtener una excepción.

\item \text{[0.25 puntos]} Al inicializar un \texttt{Verifier} con valores primos de $p$ y $q$, y con un elemento $g \in \mathbb{Z}_p^*$ tal que $|\langle g \rangle| \neq q$, se debe obtener una excepción.

\item \text{[0.5 puntos]} Para un conjunto válido de $\ell$ participantes inicializados con los mismos parámetros $p$, $q$ y $g$, si uno de ellos genera una firma para una lista de claves públicas \texttt{public\_keys}, un tag \texttt{tag} y un mensaje \texttt{message}, entonces la firma debe tener la forma $([s_1,\ldots,s_\ell], c_1, f)$ y el método \texttt{verify\_signature} con estos parámetros debe retornar \texttt{(True, True)}.


\item \text{[0.25 puntos]} Para un conjunto válido de $\ell$ participantes inicializados con los mismos parámetros $p$, $q$ y $g$, si un mismo participante genera dos firmas válidas para mensajes distintos usando la misma lista de claves públicas \texttt{public\_keys} y el mismo tag \texttt{tag}, entonces después de verificar la primera firma, el verificador debe retornar \texttt{(True, False)} al verificar la segunda firma.

\item \text{[0.25 puntos]} Para una firma válida $\sigma = ([s_1,\ldots,s_\ell], c_1, f)$ para una lista de claves públicas \texttt{public\_keys}, un tag \texttt{tag} y un mensaje \texttt{message}, si se define una firma $\sigma' = ([s'_1,s_2,\ldots,s_\ell], c_1, f)$ con $s'_1 \not\equiv_q s_1$, entonces \texttt{verify\_signature(public\_keys, message, tag, $\sigma'$)} debe rechazar la firma.
\end{itemize}
